\chapter{Kết luận}
\label{Chapter5}
Nhằm giải quyết các câu hỏi nghiên cứu về hiện tượng chuyển giao tiêu cực (negative transfer) trong quá trình thích ứng nhanh của \acrfull{maml} đã được phân tích tại Mục~\ref{sec:research_urgency}, luận văn đã đề xuất thành công phương pháp hậu diễn giải (post-hoc explaination) dựa trên độ lợi thích ứng $\Delta M$ và phương pháp \acrfull{fama}.

Cụ thể, thông qua đại lượng $\Delta M$, chiều hướng tác động (có lợi hoặc có hại) của việc thích ứng nhanh hoàn toàn có thể được xác định một cách tường minh. Mà dựa trên kết quả đó, \acrshort{fama} thông qua việc phân bổ liên hợp trong không gian đặc trưng chung $F$ mà định vị chính xác các vùng không gian trên tập hỗ trợ $S_i$ là nguyên nhân gốc rễ gây ra tác động này xuyên suốt toàn bộ quỹ đạo thích ứng qua nhiều bước, nhờ đó mà có thể trực quan hóa được hiện tượng "thiên lệch ngữ cảnh" (contextual bias) dẫn đến hiện tượng chuyển giao tiêu cực.

Đồng thời, các thực nghiệm kiểm chứng cũng đã khẳng định phương pháp mà chúng tôi đề ra hoàn toàn trung thực và đáng tin cậy với việc chỉ ra cơ chế thích ứng nhanh. Các kết quả từ việc thực hiện \acrfull{biadt} đạt mức ý nghĩa thống kê cao ($p<0.001$) ở cả hai chiều tích cực lẫn tiêu cực. Và, hai bài kiểm tra tính tỉnh táo (sanity checks) qua can thiệp tri thức tiên nghiệm $\theta_0$ và dữ liệu $S_i$ đã chứng minh bản đồ quy kết của \acrshort{fama} thực sự bắt nguồn từ tri thức đã học và nội dung dữ liệu của tập hỗ trợ, điều này đồng nghĩa với việc triệt tiêu hoàn toàn sự chi phối của các yếu tố cấu trúc mạng trơ. Nhờ đó, luận văn mang đến một công cụ diễn giải hoàn toàn mới, khả vi và độc lập với dữ liệu meta-training cho lĩnh vực meta-learning.

\section{Hạn chế và hướng phát triển}

Tuy nhiên, nghiên cứu hiện tại vẫn còn một số hạn chế do việc sử dụng kiến trúc nông. Việc thực nghiệm trên kiến trúc Conv4 khiến lưới đặc trưng ở khối tích chập cuối có độ phân giải không gian thô, gây ra hiện tượng "tràn màu" (spatial bleeding) khi định vị các vật thể nhỏ và tạo ra mức sai số biên khá lớn trong đánh giá \acrshort{biadt}. Ngoài ra, kiến trúc nông cũng khiến kịch bản can thiệp OOD chưa phản ánh đúng độ nhạy kỳ vọng do sự tương đồng về thống kê thị giác cấp thấp giữa các tập dữ liệu, và thang đo phá hủy mạng ($\ell=0,\dots,4$) chưa đủ mịn để quan sát chi tiết sự sụt giảm hiệu năng.

Từ những hạn chế trên, luận văn đề xuất 4 hướng phát triển tiềm năng trong tương lai:
\begin{enumerate}
    \item \textbf{Cải thiện độ phân giải không gian:} Mở rộng thực nghiệm sang các kiến trúc backbone sâu hơn (như ResNet-12) nhằm giải quyết triệt để hiện tượng tràn màu.
    \item \textbf{Tinh chỉnh kịch bản can thiệp:} Sử dụng các tập dữ liệu phi-vật-thể (như Places hoặc ảnh kết cấu) làm nhiễu OOD để tạo ra phép thử Sanity Check khắc nghiệt và chính xác hơn.
    \item \textbf{Kiểm chứng tính phổ quát:} Khái quát hóa khung đánh giá $\Delta M$ và \acrshort{fama} sang các biến thể tối ưu hóa khác (ANIL, BOIL, Meta-SGD) trong số các phương pháp thuộc họ \acrshort{maml} (\acrshort{maml}-based).
    \item \textbf{Can thiệp chủ động:} Kết hợp \acrshort{fama} với các cơ chế can thiệp biểu diễn (như HSFM) nhằm tiến xa hơn từ việc chẩn đoán sang việc tự động hiệu chỉnh tập hỗ trợ $S_i$, khép kín vòng lặp giữa diễn giải XAI và cải thiện trực tiếp hiệu năng mô hình.
\end{enumerate}
